

3.17 控制算法主流程
\begin{figure}[htbp]
  \centering
  \includegraphics[width=0.78\textwidth]{pic/fig3-2.png}
  \caption{全过程算法流程图}
\end{figure}








3.18 本章小结
本章将原有飞机辅助感知与自动接近控制内容整合为同一技术链路。首
先，围绕飞机前起落架对接区域设计了目标标识、视觉识别和近距测量方案，
并给出了目标身份确认、相对位姿误差和安全边界判定方法；其次，在第二章
双桥转向运动学模型基础上，建立了车体系误差动力学模型，设计了预对接
点、对接点、速度调度、误差反馈控制和四轮转向映射方法；最后，通过终端
停车集合与模式切换规则，将目标确认、接近、姿态对准、微动对接和停车保
持组织为完整闭环。
与单纯的路径跟踪控制相比，本章方法更加贴合机场低速高精度对接工
况，能够把目标身份确认、局部态势感知、轨迹可行性、误差收敛性和停车安
全性放在同一控制框架中，为后续仿真实验与系统实现提供算法基础。










\section{仿真系统构建与实现}
\subsection{系统总体架构与软硬件环境}
本文仿真系统按“操作系统—ROS 通信—Gazebo 物理环境—控制接口—自
定义算法节点”组织。Ubuntu 18.04 提供运行环境，ROS Noetic 负责话题、参
数和 TF 关系，Gazebo 11 负责刚体、碰撞和传感器仿真；gazebo\_ros\_pkgs 与
gazebo\_ros\_control 把模型、传感器和关节控制接入 ROS。自定义节点主要完
成飞机辅助感知、局部目标生成、对接控制和任务调度[11-13,16]。在调试
中，关节是否能被 controller\_manager 正确加载，是车辆能否稳定运动的关
键。
工程目录按功能拆分为四类包：tug\_description 保存 URDF/xacro、STL 网
格和材料文件；tug\_gazebo 保存 world、spawn 脚本和传感器插件；
tug\_control 保存控制器参数、话题接口和控制脚本；tug\_bringup 保存联调用
launch 文件。这样布置后，修改模型、换场景或调控制参数时不需要在同一个
文件夹里反复查找。
Gazebo 官方示例通常建议将机器人描述、仿真世界和启动文件分开管理
[11-16]。本文沿用这一思路：牵引车本体放在 description 包中，机场地面和
飞机目标放在 world 文件中，控制器加载和测试流程由 bringup 统一启动。该
方式更适合后续替换 STL 模型或增加传感器。
下表给出本文采用的基准软件栈。（表 4-1）
表 4-1 基准软件栈

组件             版本/状态                用途

Ubuntu          18.04 LTS         主机操作系统
ROS         Noetic Ninjemys    节点通信、参数、消
息、TF
Gazebo              11          刚体动力学、传感器
与场景仿真
gazebo\_ros\_pkg   Noetic 对应发行        Gazebo 与 ROS 接口
s                                    桥
gazebo\_ros\_con   Noetic 对应发行        ros\_control 接入
trol                               Gazebo
ros\_control     ROS 1 经典控制框架       控制器抽象与管理







robot\_state\_pu    Noetic 对应发行            根据 URDF
blisher                           + joint\_states 发布
TF
tf2\_ros        Noetic 对应发行          坐标变换管理
xacro         Noetic 对应发行          宏化 URDF 生成
image\_pipeline    Noetic 对应发行        相机标定与图像处理
链
OpenCV       ROS Noetic 兼容版本     飞机辅助感知标识识
别
系统架构图用于说明 ROS 与 Gazebo 之间的数据流：Gazebo 提供模型状态和
传感器数据，ROS 节点完成误差计算、状态切换和控制命令发布，ros\_control
再把命令落到转向关节和车轮关节上。
\subsection{牵引车模型与传感器建模}
URDF 采用 link-joint 树结构描述牵引车本体。为了让模型能够在 Gazebo
中参与物理仿真，仅有 visual 网格是不够的，还需要给关键部件配置
collision、inertial、Gazebo 标签和 transmission[13-16]。机场地面、飞机
目标和障碍物不放入车体 URDF，而是放在 world 文件中，使车辆模型保持清晰
的树形结构[12-14]。
本文的牵引车模型按控制需求拆分：车体、前转向桥、后转向桥、四个驱
动车轮以及前向相机、深度相机、激光雷达分别作为独立 link 或固定安装部件
处理。关节层面包括前后桥 revolute 转向关节、四个 continuous 车轮关节和
传感器 fixed 关节。这样的拆分便于分别调试转向位置控制、轮速控制和传感
器数据输出[14,16]。
车辆模型由 SolidWorks 建立。考虑到虚拟机运行 Gazebo 时算力有限，本
文没有直接使用完整高细节机械模型，而是准备了两套模型：仿真模型保留车
体外形、轮组、转向桥和传感器安装位置，删除对运动控制影响较小的细节；
机械模型面向后续样机，需要继续考虑结构强度、电气布置和维护空间。正文
中仅展示用于仿真的简化模型，如图 4-1 所示。











\begin{figure}[htbp]
  \centering
  \includegraphics[width=0.86\textwidth]{pic/fig4-1.png}
  \caption{模型图总览}
\end{figure}


模型导入 Gazebo 前，需要先确定各运动部件的关节类型。车轮采用
continuous 关节表示持续转动，传感器和附件采用 fixed 关节固定在车体上，
前后转向桥采用 revolute 关节表示有限角度转向。表 4-2 给出一组用于复现仿
真的 URDF 结构定义，后续若取得实测尺寸、质量或 CAD/BOM 参数，可直接替换
表中对应项。
表 4-2 车辆主要 URDF 结构定义

组件         类型           命名            说明

车体根坐标        link      base\_link    仅作平面参考，不带惯
量
主车体         link      base\_link    车体惯量、碰撞主体
前桥         link     front\_steer    前转向桥中间体
\_link
后桥         link     back\_steer\_    后转向桥中间体
link
前桥转向      revolute   front\_steer     绕 z 轴转向
后桥转向      revolute    rear\_steer     绕 z 轴转向
相机        fixed     front\_camer   识别飞机辅助码主相机
a\_joint









深度相机                   fixed     depth\_camer    近距深度感知
a\_joint
激光雷达                   fixed     lidar\_joint   前向近距安全监测
在此示例部分关节定义：
<link
name="base\_link">
<inertial>
<origin

xyz="0.052294 1.5562E-06 0.061913"
rpy="0 0 0" />
<mass

value="34.871" />
<inertia
ixx="0.51812"
ixy="-2.2269E-05"
ixz="0.18429"
iyy="2.6761"
iyz="-2.8808E-06"
izz="3.0257" />
</inertial>

<visual>
<origin
xyz="0 0 0"

rpy="0 0 0" />
<geometry>
<mesh

filename="package://zongzhuang/meshes/base\_link.STL" />
</geometry>
<material

name="">
<color







rgba="0.79216 0.81961 0.93333 1" />
</material>
</visual>
<collision>
<origin
xyz="0 0 0"
rpy="0 0 0" />
<geometry>

<mesh
filename="package://zongzhuang/meshes/base\_link.STL" />
</geometry>

</collision>
</link>
\subsection{任务与模式管理}
任务与模式管理器相当于控制链路中的上层状态机。它不直接求解前后桥转
角和车轮速度，而是根据目标身份、相对位姿、障碍物距离和车辆当前状态，
决定系统处于哪个作业阶段，并向下层控制器给出目标点类型、速度上限和是
否允许继续推进等标志量[5,7,25,28]。本文采用以下几类模式：
搜索模式：目标身份尚未确认时，车辆不进入接近流程，轮速保持为零，转
向桥回中或低速等待识别结果。
自动接近模式：目标身份有效且距离较远时，车辆以预对接点为目标，允许
较高速度上限和较积极的转向修正。
姿态对准模式：车辆进入近距范围后，速度上限降低，控制重点转为收敛横
向误差和航向误差。
微动对接模式：最终接近阶段只保留小幅速度和转向输出，使牵引车缓慢靠
近对接点。
停车模式：对接条件满足或出现异常中止时，系统发布零轮速命令，并使前
后桥转向回中。




\subsection{系统全流程测试}







为检验系统能否连续完成作业链路，本文在 ROS 与 Gazebo 中设计了自动牵引
车全流程测试。测试不是单看某一次行驶动作，而是按任务顺序依次执行车辆
前往飞行器区域、对接飞行器、推出飞行器、解除牵引和车辆撤离五个阶段
[5-7]。通过该流程，可以检查车辆速度、姿态、转向命令和状态机切换是否
保持一致。
\paragraph{} 1. 车辆前往飞行器区域
测试开始前，先将牵引车放置在模拟航站楼附近，并让转向桥完成回中、
轮速归零和短暂停车保持，以排除模型刚生成时的轻微滑动影响。随后车辆沿
地勤车辆通道驶向目标飞行器区域。该阶段主要检查三个动作是否连贯：车辆
能否按指令起步，轮速命令能否稳定驱动车体前进，前后桥转向是否能在道路
方向变化时给出可观察的姿态修正。
在这一段仿真中，速度没有设置得很高，主要是为了观察车辆在长距离接
近时是否会出现摆头、跑偏或转向回中不及时。车辆接近飞行器区域后，重点
查看车身朝向、路径连续性以及与预对接区域的相对位置；若此处已经出现明
显偏差，后续对接阶段即使继续执行也容易放大误差。因此，图 4-2 主要用于
说明牵引车能够稳定进入对接准备区，而不是单纯展示车辆移动动画。




\begin{figure}[htbp]
  \centering
  \includegraphics[width=0.86\textwidth]{pic/fig4-2.png}
  \caption{车辆在航站楼下运行，前往目标飞行器}
\end{figure}


\paragraph{} 2. 车辆对接飞行器
牵引车进入飞行器附近后，系统切换到对接控制阶段。此时车辆不再只是
沿道路行驶，而是让牵引点逐步靠近飞行器牵引点，并在位置和姿态满足要求
时完成接近。控制节点根据感知模块给出的相对位姿信息，计算距离误差、横
向误差和航向误差，再修正车速与转向角。（图 4-3）
对接过程中，车辆先到达预对接位置，再以低速驶向最终对接点。与普通
路径跟踪相比，该阶段更关注末端稳定性，因此需要限制速度和转向变化率，








避免车身摆动过大或对接误差继续放大。（图 4-4）




\begin{figure}[htbp]
  \centering
  \includegraphics[width=0.86\textwidth]{pic/fig4-3.png}
  \caption{尝试对接。此时车辆没有正对飞行器，正在通过蟹行平移进行调整}
\end{figure}





\begin{figure}[htbp]
  \centering
  \includegraphics[width=0.86\textwidth]{pic/fig4-4.png}
  \caption{对接成功，拖把末端固定装置放下。等待推出命令}
\end{figure}


该阶段用于检验牵引车能否依据飞行器目标位置完成姿态调整，并在低速
状态下稳定接近指定对接点。
\paragraph{} 3. 车辆推出飞行器
完成对接后，系统进入推出阶段。此时牵引车与飞行器在仿真中按整体运
动关系处理，牵引车沿预设推出方向把飞行器从停放区域推至目标位置。该阶
段重点观察速度是否平稳、转向是否连续，以及飞行器姿态是否出现突变。
推出过程中，控制系统主要检查车辆速度曲线、转向命令和飞行器整体轨
迹。由于牵引车与飞行器之间存在相对约束，相比单车行驶更容易暴露路径偏
差、牵引姿态摆动或转向不连续等问题。（图 4-5—图 4-7）











\begin{figure}[htbp]
  \centering
  \includegraphics[width=0.86\textwidth]{pic/fig4-5.png}
  \caption{飞行器被推出过程}
\end{figure}





\begin{figure}[htbp]
  \centering
  \includegraphics[width=0.86\textwidth]{pic/fig4-6.png}
  \caption{拖把解除，车辆倒退一部分距离等待飞行器报告}
\end{figure}





\begin{figure}[htbp]
  \centering
  \includegraphics[width=0.86\textwidth]{pic/fig4-7.png}
  \caption{车辆撤离}
\end{figure}


\section{解除牵引}
飞行器推出到指定位置后，牵引车先减速停车，再执行解除牵引动作。连
接关系解除后，飞行器保持在目标位置，牵引车重新进入独立运动状态。该阶
段用于检查“停止—解挂—后退”过程是否连贯。
解除牵引时，重点观察牵引车是否停稳、飞行器是否产生额外位移，以及
两者之间是否保留安全间隔。只有这些条件满足后，系统才允许车辆进入撤离







阶段。（图 4-6）





\paragraph{} 5. 车辆撤离
完成解除牵引后，牵引车按照预设撤离路径离开飞行器区域。撤离阶段主
要验证牵引车在完成作业任务后能否安全驶离，并避免与飞行器或周围环境发
生干涉。系统根据撤离方向控制车辆转向与速度，使其逐步远离飞行器，并回
到指定待命区域或安全区域。
在该阶段，车辆应保持稳定行驶状态，转向过程应尽量平缓，避免因急转
或急停影响仿真效果。车辆撤离后，整个自动牵引对接与推出流程结束。（图
4-7）





\subsection{本文方法与人工操作以及其它方法对比分析}
为进一步说明本文所设计系统的特点，有必要将其与传统人工牵引方式、
固定路径控制方式以及常见自动驾驶控制算法进行对比分析。
\subsubsection{与传统人工牵引方式的比较}
传统机场牵引主要依靠驾驶员、地面指挥员和现场协同人员完成。人工方式的






优势在于现场判断灵活，能够根据飞机位置、障碍物和临时指令调整操作；不
足是对人员经验和沟通质量依赖较强，重复作业精度也难以完全一致。
本文所提出的方法主要面向仿真环境下的自动化流程验证。相比人工牵引
方式，本文方法的优势并不在于真实复杂环境中的经验判断能力，而在于能够
将牵引作业过程转化为可建模、可重复、可分析的控制流程。通过 ROS 节点通
信、Gazebo 物理仿真、车辆运动控制和目标位姿误差反馈，可以较为清晰地表
达车辆接近、对准、推出和撤离等动作过程。
从本项目实际测试过程来看，车辆模型在 Gazebo 中能够完成基本行驶、转
向、刹停和回中等控制动作，并可通过脚本或控制节点实现分阶段作业流程。
这说明本文方法在标准化仿真验证方面具有一定价值。下图给出了本文方法与
其它方法的比较柱状图（图 4-8）。




\begin{figure}[htbp]
  \centering
  \includegraphics[width=0.78\textwidth]{pic/fig4-8.png}
  \caption{本文方法与主流方法以及理论方法比较}
\end{figure}


\subsubsection{与固定路径或开环脚本控制方式的比较}
相比与其它方式而言，最简单实现机场自动牵引驳接的方式是采用固定路径
或开环脚本控制，即预先设定车辆的速度、转向角和运动时间，使车辆按照固定
流程完成前进、转弯、倒车或停止等动作。但固定路径或开环控制的局限性也比
较明显。由于其控制过程主要依赖预设时间和固定参数，当车辆初始位置、飞机
停放位置或仿真环境发生变化时，车辆无法根据实际偏差自动修正运动状态。例
如，在本项目调试过程中，车辆曾出现初始缓慢滑动、转向不完全回中、不同速
度下运动效果不一致等现象，如果完全依靠固定脚本控制，容易造成路径偏差累
积，影响后续对接效果。






本文方法在固定流程的基础上引入了相对位姿误差思想，将车辆与目标对
接点之间的位置误差、横向误差和航向误差作为控制依据，并结合预对接点、
最终对接点和分阶段模式切换来组织整个作业流程。因此，相较于单纯开环脚
本，本文方法具有更强的闭环控制特征，能够更合理地描述自动牵引车在接近
目标过程中的修正行为（图 4-9）。




\begin{figure}[htbp]
  \centering
  \includegraphics[width=0.78\textwidth]{pic/fig4-9.png}
  \caption{几种方法对接过程中的横向误差收敛对比}
\end{figure}


\subsubsection{与常见路径跟踪算法的比较}
在移动机器人和智能车辆研究中，常见路径跟踪方法包括 PID 控制、Pure
Pursuit 算法、Stanley 算法以及模型预测控制等。
PID 控制结构简单、实现方便，适合速度控制、角度控制等单一变量调节。
在本项目中，车辆的车轮速度控制器和转向关节控制器本身就需要依赖底层控
制参数实现稳定执行。但如果仅采用简单 PID 完成整体对接任务，则难以同时
处理车辆纵向距离、横向偏差和航向角误差之间的耦合关系。牵引车对接飞行
器时不仅要“走到目标附近”，还要保证车身姿态与牵引点方向基本一致，因
此单一 PID 控制难以完整表达对接过程。
Pure Pursuit 适合一般路径跟踪，可根据前视点生成较平滑的转向控制
量。但近距离对接时，车辆与目标点距离短、速度低，末端姿态要求高，单纯
依赖前视点容易出现最后几厘米或最后几度角对不准的问题。本文方法更强调
车辆与目标之间的相对位姿，因此更适合低速、近距、强姿态约束的对接场
景。
Stanley 算法常用于自动驾驶路径跟踪，特别适合处理中高速行驶中的横向
误差。但机场牵引车对接飞行器属于低速精细作业，控制重点不是高速循迹，
而是小范围内的位置修正和姿态调整。因此，本文采用相对误差和阶段切换结
合的方式，而不是直接套用 Stanley 方法。








MPC 能够同时考虑车辆模型、转向角限制、速度限制和未来轨迹优化，理论
上适合复杂约束控制。但 MPC 对模型精度、求解器实时性和参数调试要求较
高。结合本项目时间和实现条件，直接采用 MPC 会显著增加系统复杂度。因
此，本文选择相对位姿误差控制和规则化流程管理，在可实现性与任务完整性